\paragraph{Frozen baseline supplement.} N1024 comparison on the original common500 sources and, separately, the previously published new100 sources. Every point averages three registered noise replicates per source. Shading shows pointwise 95\% source bootstrap intervals. Original method intervals are reused unchanged; only new method intervals and required source-paired contrasts are computed with the original 10,000-replicate implementation and seed 2026100701. These are subsequent baseline comparisons on existing evaluation cohorts, not newly blind tests. All failed frames remain in the denominators.

\paragraph{Paired contrasts.} Deltas are method minus reference after averaging each source over its three noise replicates; error bars are the directly computed source-paired 95\% intervals. LPIPS retains its original sign, so negative deltas are better. ConvNeXt agreement and its intervals are shown as percentages in the main curves and percentage points in contrasts, not relative gains or classification accuracy. The common500 and new100 populations are never combined or subtracted from one another.

\paragraph{SwinJSCC qualification.} The 80k adapted baseline at 19 dB is outside its training and calibration range, shown with an open marker and a dashed final segment. The continued adapted model uses the authorized weights-only initialization with fresh optimizer, sampling and random states, and expanded 1--19 dB training coverage. It reaches its 120k additional-update budget while still improving; neither official optimality nor convergence is claimed. P1024 resumes its full original training state and meets the registered empirical plateau rule at 100k.
